% Standalone problem document, generated from the adjacent README.md.
% Compile with XeLaTeX. Mathematical content is maintained in Markdown.
\documentclass[11pt,a4paper]{article}
\usepackage[fontsize=10.5pt]{scrextend}
\usepackage[margin=22mm,headheight=15pt,headsep=9mm,footskip=11mm]{geometry}
\usepackage{fontspec}
\setmainfont{texgyrepagella}[Extension=.otf,UprightFont=*-regular,BoldFont=*-bold,ItalicFont=*-italic,BoldItalicFont=*-bolditalic]
\setsansfont{texgyreheros}[Extension=.otf,UprightFont=*-regular,BoldFont=*-bold,ItalicFont=*-italic,BoldItalicFont=*-bolditalic]
\setmonofont{texgyrecursor}[Extension=.otf,UprightFont=*-regular,BoldFont=*-bold,ItalicFont=*-italic,BoldItalicFont=*-bolditalic,Scale=MatchLowercase]
\usepackage{amsmath,amssymb}
\usepackage{unicode-math}
\setmathfont{texgyrepagella-math.otf}
\usepackage{xcolor,microtype,enumitem,needspace,fancyhdr}
\usepackage{bookmark}
\definecolor{ink}{HTML}{17344B}
\definecolor{accent}{HTML}{197D80}
\definecolor{muted}{HTML}{596773}
\hypersetup{colorlinks=true,linkcolor=accent,urlcolor=accent,pdftitle={SP-06 - A
real-valued symbol on a Jordan curve and real Toeplitz
spectra},pdfauthor={Open Problems in Numerical Linear Algebra}}
\urlstyle{same}
\setlength{\parindent}{0pt}
\setlength{\parskip}{5pt plus 1pt minus 1pt}
\linespread{1.04}
\setlength{\emergencystretch}{3em}
\widowpenalty=10000
\clubpenalty=10000
\displaywidowpenalty=10000
\allowdisplaybreaks[1]
\setlist{nosep,leftmargin=1.5em}
\providecommand{\tightlist}{\setlength{\itemsep}{0pt}\setlength{\parskip}{0pt}}
\providecommand{\pandocbounded}[1]{#1}
\pagestyle{fancy}
\fancyhf{}
\fancyhead[L]{\sffamily\footnotesize\color{muted}OPEN PROBLEMS IN NUMERICAL LINEAR ALGEBRA}
\fancyhead[R]{\sffamily\bfseries\footnotesize\color{accent}SP-06}
\fancyfoot[L]{\parbox[b]{0.9\textwidth}{\sffamily\fontsize{8}{10}\selectfont\color{muted}Editorial ratings; a bounded search cannot certify that a problem remains unsolved.\\Literature check: 2026-09-08}}
\fancyfoot[R]{\sffamily\footnotesize\color{muted}\thepage}
\renewcommand{\headrulewidth}{0.3pt}
\renewcommand{\headrule}{\hbox to\headwidth{\color{accent}\leaders\hrule height \headrulewidth\hfill}}
\makeatletter
\renewcommand{\section}{\Needspace{5\baselineskip}\@startsection{section}{1}{0pt}{12pt plus 2pt minus 2pt}{4pt}{\sffamily\bfseries\large\color{ink}}}
\renewcommand{\subsection}{\Needspace{4\baselineskip}\@startsection{subsection}{2}{0pt}{9pt plus 1pt minus 1pt}{3pt}{\sffamily\bfseries\normalsize\color{ink}}}
\makeatother
\setcounter{secnumdepth}{0}
\begin{document}
{\sffamily\small\color{muted}Eigenvalues and inverse problems\par}
\vspace{3pt}
{\raggedright\hyphenpenalty=10000\exhyphenpenalty=10000\sffamily\bfseries\fontsize{21}{24}\selectfont\color{ink}A
real-valued symbol on a Jordan curve and real Toeplitz spectra\par}
\vspace{7pt}
{\sffamily\small\textbf{Difficulty:} challenging\quad\textbf{Importance:} interesting
to the community\par}
\vspace{4pt}
{\color{accent}\hrule height 0.8pt}
\vspace{6pt}
\textbf{Topic:} Spectra of finite banded Toeplitz matrices.

\subsection{Problem statement}\label{problem-statement}

Let \(r,s\ge1\) be integers and

\[
b(z)=\sum_{k=-r}^{s}b_kz^k,
\qquad b_k\in\mathbb C,\qquad b_{-r}b_s\ne0.
\]

For every \(n\ge1\), define \(T_n(b)=(b_{i-j})_{i,j=1}^n\), with
\(b_k=0\) outside \([-r,s]\). Suppose there is a Jordan curve
\(\gamma\subset\mathbb C\setminus\{0\}\) such that \(b(z)\in\mathbb R\)
for every \(z\in\gamma\). A Jordan curve means the image of a continuous
injective map from the unit circle.

Must

\[
\operatorname{spec}(T_n(b))\subset\mathbb R
\qquad\text{for every integer }n\ge1?
\]

\subsection{Why it matters}\label{why-it-matters}

The conjecture would characterize a broad source of real spectra in
nonsymmetric structured eigenvalue problems through a scalar symbol
condition, uniformly over matrix size.

\subsection{References}\label{references}

\begin{itemize}
\tightlist
\item
  B. Shapiro and F. Štampach,
  \href{https://arxiv.org/abs/1702.00741}{Non-Self-Adjoint Toeplitz
  Matrices Whose Principal Submatrices Have Real Spectrum},
  \emph{Constructive Approximation} 49 (2019), 191--226. Equation (2),
  Theorem 1(ii)--(iii), and Theorem 8; use arXiv v4, which includes the
  correction.
\item
  B. Shapiro and F. Štampach,
  \href{https://doi.org/10.1007/s00365-022-09614-0}{Correction to the
  same article}, 2023. The correction, pp.~27--28 of the combined arXiv
  PDF, withdraws the proof of (ii)\(\Rightarrow\)(iii) and explicitly
  proposes the implication as conjectural.
\item
  D. Giandinoto, \href{https://arxiv.org/abs/2411.16266}{On reality of
  eigenvalues of banded block Toeplitz matrices}, 2024, introduction and
  §2. The paper describes the corrected scalar implication as a
  conjecture before considering a block generalization.
\end{itemize}

\subsection{Status check}\label{status-check}

On 2026-09-08, checked the corrected arXiv v4 (2022-12-29), its appended
erratum, and Giandinoto's November 2024 paper. Searches combined
``Shapiro'', ``Štampach'', ``Toeplitz'', ``reality'', ``erratum'',
``Jordan curve'', ``proof'', and 2026. No general resolution was
located. The original article's theorem label is misleading without its
correction: the finite-spectrum implication above remains the authors'
stronger proposed statement. Its weaker limiting-spectrum consequence
and particular symbol families are not separate catalog problems here.
\end{document}
